% Introduction (English draft v1, LaTeX)
% Requires in preamble:
%   \usepackage{xspace}
%   \newcommand{\sys}{HyperFront\xspace}
% Citation keys are placeholders; replace with your .bib entries.

\section{Introduction}
\label{sec:intro}

In recent years, driven by the rapid growth of AI, RDMA has been widely
deployed to serve AI/ML applications for its high bandwidth, low latency,
and low CPU overhead. These performance benefits are now prompting cloud
providers to bring RDMA into their public cloud networks
\cite{azure-rdma,alibaba-erdma}.

Compared with building a dedicated RDMA network, deploying RDMA inside
the existing VPC network has obvious appeal: it adds no extra equipment
or management overhead, and it retains full flexibility for the diverse
needs of tenant applications. The price, however, is that RDMA traffic
must compete for bandwidth and other network resources with the TCP
traffic that already lives in the VPC---a competition RDMA is poorly
equipped for. RDMA was originally designed for single-tenant
environments with homogeneous traffic (e.g., HPC), typically assuming
shallow-buffered fabrics where packet loss is rare; in a multi-tenant
VPC carrying complex traffic mixes, RDMA turns out to be remarkably
fragile. Hardware Traffic Classes (HW TCs) do support traffic isolation,
but their scarcity means they can at best separate traffic
\emph{types}, falling short of the dual goal of isolating both tenants
and traffic classes; moreover, their static bandwidth weights are a poor
match for the complex and rapidly shifting traffic patterns of a VPC.

Unfair bandwidth competition with TCP arises both at the end hosts and
at the switches. At the sending side, the bottleneck's owner and all of
its competitors reside on the same host, so local scheduling suffices to
resolve the contention. The genuinely hard cases are the last-hop switch
and the receiving end, where the competing flows originate from many
different machines.

When incast occurs, TCP and RDMA compete for the egress bandwidth of the
last-hop switch, and the heterogeneous congestion control of the two
classes amplifies the unfairness. Take RDMA-DCQCN~\cite{dcqcn} and
TCP-Cubic~\cite{cubic} as an example: switches typically start marking
ECN at a queue depth well below the egress buffer limit, triggering
DCQCN's rate reduction, whereas Cubic backs off only upon packet loss.
RDMA thus yields far earlier than TCP, and after convergence RDMA is
left with a very small share of the bandwidth.

The bottleneck can also sit at the receiving VM: cloud providers sell
VMs with a downlink bandwidth cap below the physical line rate. The
common receiver-side policing practice is simply dropping excess
packets. This bottleneck is peculiar: from the network's perspective
there is no congestion at all---the scarcity is contractual rather than
physical---so the policer emits none of the congestion signals that
DCQCN depends on (in our measurements, senders received zero CNPs;
\S\ref{sec:motivation}). RDMA is then left with only loss recovery to
fall back on, and even lossy RDMA is far more sensitive to packet loss
than TCP, so RDMA again converges to a very low bandwidth. ByteDance's
JAKIRO~\cite{jakiro} proposes DHTB to address this problem, bringing
receiver-side congestion back into the scope of congestion control. Yet,
owing to the same CC heterogeneity, the fairness JAKIRO provides is
neither accurate nor stable: in our experiments, the very same
configuration randomly converges either to the fair share or to a
runaway state in which RDMA grabs 2--3$\times$ the purchased cap, and
whether it converges at all is dictated by how aggressive the sender's
DCQCN recovery parameters happen to be (\S\ref{sec:motivation}). In
other words, JAKIRO manufactures the missing signal, but outsources
fairness to the cooperation of the tenants' congestion control.

Putting the two venues of competition together, the crux of the
unfairness problem is that each protocol's design is blind to the other
class. When the two classes of traffic run side by side, the bandwidth
each ends up with is merely the by-product of independently designed
congestion-control algorithms colliding with each other---shaped by
algorithm design and parameters, network queueing, and more---and the
convergence point is unpredictable. EQDS~\cite{eqds} recognizes this as
well, and therefore wraps all traffic into a unified datagram layer,
moving queueing to the sender and having the receiver issue credits so
that all traffic is scheduled uniformly at the end hosts. However, EQDS
is too heavy to deploy in a VPC network to be realistic; and for
fairness as the sole objective, its full design is far more than what is
needed.

Reflecting on the limitations of existing solutions, we argue that a
good solution must satisfy the following requirements:

\begin{itemize}
  \item \textbf{[R1] Inter-class fairness with dual isolation:} resolve
  the bandwidth fairness problem between RDMA and TCP in cloud networks,
  providing both tenant-level and class-level isolation.
  \item \textbf{[R2] (Scalability) High-performance and accurate:} must
  not become a bottleneck on the data plane under large numbers of
  concurrent flows, while enforcing rate limits accurately and stably.
  \item \textbf{[R3] Deployment ready:} transparent to tenants and
  applications---no modification to application APIs---and compatible
  with commodity NICs and switches; in particular, for RDMA, it must
  support diverse RDMA operations.
\end{itemize}

To this end, we present \sys, a solution focused on TCP/RDMA fairness
that is deployable in VPC networks. Our key insight is that although the
competition may arise either in the network or at the end host, the
receiving end has sufficient visibility to distinguish the two, and can
characterize the overall state of both kinds of competition---accurately
and stably---through one unified signal. The system consists of three
modules: the Traffic Shaper, the Traffic Meter, and the TX/RX Agents.
The Traffic Meter runs on the receiving end and continuously measures
the arrival rate of each \emph{flow-set} (all traffic of one class from
one source VM to one destination VM), metering TCP and RDMA separately;
the Traffic Shaper enforces rate limits at the sending end---driving the
NIC's hardware scheduler for RDMA and kernel EDT pacing for TCP---both
sitting in the infrastructure domain, invisible to tenant applications.
On top of them, the RX Agent uses the Meter's data to quantify the
degree of unfairness along the entire network path and feeds this signal
back to the TX Agent; the TX Agent underpins all protocols with a
single, CC-agnostic MIMD rate-control law that caps the sending rate of
each flow-set and installs the caps into the Shaper.

% ------------------------------------------------------------------
% TODO (author): implementation, results, and contributions paragraphs.
% ------------------------------------------------------------------
